\thispagestyle{plain}
\phantomsection
\addcontentsline{toc}{chapter}{中文摘要}
\centerline{\zihao{-3}\heiti 中文摘要}

\linespread{1.4}\zihao{-4}\bigskip

本文基于 \texttt{CTaylorSeriesSecond20250324} 代码资料，围绕“高阶泰勒级数数值积分 + 牛顿类参数迭代”的统一计算框架开展建模分析。针对常微分方程系统的高精度求解需求，构建了状态变量的高阶导数递推表达，形成16阶局部泰勒展开积分器，并与经典Runge--Kutta方法进行精度与效率对比。

在参数估计部分，本文将模型辨识写成非线性最小二乘问题，分别给出牛顿法与高斯牛顿法的目标函数、梯度和近似海森矩阵结构；结合代码中的导数模块，设计了可复用的“函数--梯度--海森”接口。实验从步长、噪声水平、初值扰动和迭代阈值四个维度评估算法稳定性。

结果表明：高阶泰勒法在平滑动力系统中具有较高精度优势；高斯牛顿法在中等噪声下具有更好的数值稳健性；当初值较远时，阻尼策略与步长自适应可显著提升收敛成功率。本文形成了“模型推导--算法实现--数值验证”的完整闭环，可为高精度ODE求解与参数辨识任务提供参考。

\bigskip
\noindent{\zihao{4}\heiti 关键词：}
高阶泰勒级数；常微分方程；高斯牛顿法；参数辨识；数值稳定性
